\printedpage{4}
% head: A. Grothendieck : Dérivateurs, Ch. XV, Édition des Universités de Montpellier II et Paris VII
% foot: Version 5 novembre 2010

Thus, if we consider
\[
\left\{
\begin{array}{rl}
\Flu(\cM) & \text{as an object of CAT}/\cM \text{ by the functor } f \longmapsto \text{target}\, f\\[1ex]
\text{where} & \text{CAT is the category of categories \textit{elements of} } \fV\text{\textit{,} a universe}\\
 & \text{such that } \fU \in \fV\text{, whence Cat } \overset{\text{(def)}}{=} \mathrm{Cat}^{\fU},\ \cM,\ \Flu(\cM) \in \text{CAT} ,
\end{array}
\right.
\]
we have an endofunctor $p_1$ of $\Flu(\cM)$ over $\cM$, and a functorial homomorphism
\[
i_1 \ : \ \mathrm{id}_{\Flu(\cM)/\cM} \longrightarrow p_1 \ , \qquad \textit{i.e.} \quad
\begin{tikzcd}[column sep=small]
X \arrow[rr, "f"] \arrow[dr, "i_1(f)"'] & & Y \\
 & \underbrace{X_1(f)}_{=\,\text{source}\, p_1(f)} \arrow[ur, "p_1(f)"'] &
\end{tikzcd} \ ,
\tag*{(1.13)}
\]
and the description of the ordinal sequence of endofunctors $p_\alpha$ of $\Flu\,\cM/\cM$ follows from this datum alone, taking into account that small filtering limits exist in the fibers of $\Flu\,\cM \xrightarrow{\text{target}} \cM$.

\msp{page 5}

\[
\begin{tikzcd}[column sep=normal]
X \arrow[r, "i_1(f)"] \arrow[dr, "f"'] & X_1(f) \arrow[r, "i_1(p_1(f))"] \arrow[d, "p_1(f)"] & X_2(f) \arrow[r] \arrow[dl, "p_2(f)"] & \cdots \arrow[r] & X_\alpha(f) \arrow[r, "i_1(p_\alpha(f))"] \arrow[dlll, "p_\alpha(f)"] & X_{\alpha+1}(f) \arrow[r] \arrow[dllll, "p_{\alpha+1}(f) = p_1(p_\alpha(f))"] & \cdots \\
 & Y & & & & &
\end{tikzcd}
\tag*{(1.14)}
\]

So, to finish defining the ordinal sequence of the $X_\alpha(f)$ (equipped with the $i_\alpha(f)$, $p_\alpha(f)$), or of the $p_\alpha(f)$ (equipped with the $i_\alpha(f)$), it remains to describe $X_1(f)$, $i_1(f)$, $p_1(f)$.

The essential property of this triple that we shall need is the following:

\begin{itemize}[leftmargin=5em, labelwidth=4em, labelsep=1em, align=left]
\item[(1.15)] \textbf{Fact}$(C_0)$ (factorization property with respect to $C_0$).

For every commutative diagram
\[
D : \qquad
\begin{tikzcd}
A \arrow[r] \arrow[d, "i"'] & X \arrow[d, "f"] \\
B \arrow[r] & Y
\end{tikzcd}
\qquad \text{with } i \in C_0 \ ,
\]
whence a solid-arrow diagram
\[
\begin{tikzcd}
A \arrow[r] \arrow[d, "i"'] & X \arrow[d, "i_D"] \arrow[dr, dotted, "i_1(f)"] \arrow[dd, dotted, bend right=60, "f"'] & \\
B \arrow[r] & X_D \overset{\text{def}}{=} X \amalg_A B \arrow[r, dotted, "u_D"] \arrow[d, "p_D"] & X_1(f) \arrow[dl, dotted, "p_1(f)"] \\
 & Y \ , &
\end{tikzcd}
\]
\printedpage{5}
% head: A. Grothendieck : Dérivateurs, Ch. XV, Édition des Universités de Montpellier II et Paris VII
% foot: Version 5 novembre 2010
there exists a homomorphism from the factorization $f = p_D i_D$ of $f$ to the factorization $f = p_1(f)i_1(f)$.
\end{itemize}

To formulate this condition, I have used the existence of the amalgamated sums $X_D$, \textit{i.e.} the condition on $C_0$:

\msp{page 6}

\begin{itemize}[leftmargin=5em, labelwidth=4em, labelsep=1em, align=left]
\item[(1.16)] The arrows $i \in C_0$ are cosquarable, \textit{i.e.} for a diagram
\[
\begin{tikzcd}
A \arrow[r, "u"] \arrow[d, "i"'] & X \\
B &
\end{tikzcd}
\]
with $i \in C_0$, the amalgamated sum exists.
\end{itemize}

But we can also formulate 1.15 as follows, independently of the hypothesis (1.16):

\begin{itemize}[leftmargin=5em, labelwidth=4em, labelsep=1em, align=left]
\item[(1.17)] For every commutative square
\[
\begin{tikzcd}
A \arrow[r, "u"] \arrow[d, "i"'] & X \arrow[d, "f"] \\
B \arrow[r, "v"] & Y
\end{tikzcd}
\qquad ,
\]
with $i \in C_0$, there exists a solid-arrow commutative square
\[
\begin{tikzcd}
A \arrow[r, "u"] \arrow[d, "i"'] & X \arrow[d, "i_1(f)"] \\
B \arrow[r, "v_1"] \arrow[dr, dotted, "v"] & X_1(f) \arrow[d, dotted, "p_1(f)"] \\
 & Y
\end{tikzcd}
\]
such that $v = p_1(f)v_1$.
\end{itemize}

This has as a consequence the

\textbf{Corollary 1} (of 1.17). \textit{Let} $f \in \Fl(\cM)$\textit{, and} $\alpha$ \textit{be an ordinal} $(\in \fU)$. \textit{For every commutative square}
\[
\begin{tikzcd}
A \arrow[r, "u_\alpha"] \arrow[d, "i"'] & X_\alpha(f) \arrow[d, "p_\alpha(f)"] \\
B \arrow[r, "v"] & Y
\end{tikzcd}
\]
\printedpage{6}
% head: A. Grothendieck : Dérivateurs, Ch. XV, Édition des Universités de Montpellier II et Paris VII
% foot: Version 5 novembre 2010
\textit{there exists a commutative square}
\[
\begin{tikzcd}
A \arrow[r, "u_\alpha"] \arrow[d, "i"'] & X_\alpha(f) \arrow[d, "\text{trans.}"] \arrow[dd, dotted, bend left=60, "p_\alpha(f)"] \\
B \arrow[r, "v_1"] \arrow[dr, dotted, "v"] & X_{\alpha+1}(f) \arrow[d, dotted, "p_{\alpha+1}(f)"] \\
 & Y
\end{tikzcd}
\]
\textit{such that}
\[
v \ = \ p_{\alpha+1}(f)\ v_1 \ .
\]

\msp{page 7}

\textbf{Corollary 2.} \textit{Let $\pi$ be a cardinal} $(\in \fU)$ \textit{such that for every arrow $i : A \longrightarrow B$ in $C_0$, the source $A$ is $\pi$-accessible. Let $\alpha$ be an ordinal large relative to $\pi$} (i.e. \textit{such that the ordered set $I_\alpha$ of ordinals $\alpha' < \alpha$ is large relative to $\pi$}) -- \textit{e.g. $\alpha$ the first ordinal whose cardinality is $> \pi$. Then for every arrow $f \in \Fl(\cM)$, we have $p_\alpha(f) \in (C_0)_*$.}

\textsc{Proof}. This means that for every commutative square
\[
\begin{tikzcd}
A \arrow[r, "u"] \arrow[d, "i"'] & X_\alpha(f) \arrow[d, "p_\alpha(f)"] \\
B \arrow[r, "v"] \arrow[ur, dotted, "w"] & Y \ ,
\end{tikzcd}
\tag*{$(*)$}
\]
with $i \in C_0$, the dotted arrow $w$ (with $wi = u$, $p_\alpha(f)w = v$) exists. Now
\[
X_\alpha(f) = \varinjlim_{\alpha' \in I_\alpha} X_{\alpha'}(f) \ ,
\]
so $A$ being $\pi$-accessible and $I_\alpha$ large relative to $\pi$, $u$ comes from an arrow $u_{\alpha'} : A \longrightarrow X_{\alpha'}(f)$, and (if necessary replacing $\alpha'$ by $\alpha'' > \alpha'$ and $u_{\alpha'}$ by the composite $u_{\alpha''} : A \longrightarrow X_{\alpha'}(f) \longrightarrow X_{\alpha''}(f)$) we may suppose $p_{\alpha'}(f)u_{\alpha'} = vi$, \textit{i.e.} that $(*)$ comes from a commutative square
\[
\begin{tikzcd}
A \arrow[r, "u_{\alpha'}"] \arrow[d, "i"'] \arrow[rr, dotted, bend left=30, "u"] & X_{\alpha'}(f) \arrow[r, dotted, "\text{trans.}"'] \arrow[d, "p_{\alpha'}(f)"'] & X_\alpha(f) \arrow[dl, dotted, "p_\alpha(f)"] \\
B \arrow[r, "v"] & Y \ , &
\end{tikzcd}
\]

\printedpage{7}
% head: A. Grothendieck : Dérivateurs, Ch. XV, Édition des Universités de Montpellier II et Paris VII
% foot: Version 5 novembre 2010
and applying corollary 1, we find
\[
\begin{tikzcd}
A \arrow[r, "u_{\alpha'}"] \arrow[d, "i"'] & X_{\alpha'}(f) \arrow[d, "\text{trans.}"] & \\
B \arrow[r, "v_1"] \arrow[dr, dotted, "v"] & X_{\alpha'+1}(f) \arrow[r, dotted, "\text{trans.}"] \arrow[d, "p_{\alpha+1}(f)"] & X_\alpha(f) \arrow[dl, dotted, "p_\alpha(f)"] \\ % [sic] p_{\alpha+1}(f) -> p_{\alpha'+1}(f), the arrow starts at X_{\alpha'+1}(f); ?? uncertain: subscript of the label, tightly set -- \alpha+1 or \alpha'+1
 & Y &
\end{tikzcd}
\]
with $v = p_{\alpha+1}(f)v_1$, and we shall take % [sic] p_{\alpha+1}(f)v_1 -> p_{\alpha'+1}(f)v_1, corollaire 1 is applied to \alpha' (arrow from X_{\alpha'+1}(f))
\[
w \ = \ \mathrm{trans}_{\alpha,\alpha'+1} \circ v_1 \ .
\]

\msp{page 8}

But we want a factorization of $f$, an arbitrary arrow in $\cM$, as
\[
f \ = \ p\ i \ , \qquad \text{with} \ \ p \in (C_0)_* \ \ \textit{et} \ \ i \in \widetilde{C_0} \ ,
\]
so we must still state a property, ensuring that the $i_\alpha(f)$ are in $\widetilde{C_0}$. We shall even have something more precise. Let us introduce

\begin{itemize}[leftmargin=10em, labelwidth=9em, labelsep=1em, align=left]
\item[(1.18)\hfill $C'_0 \ = $] smallest subset of $\Fl(\cM)$ containing $C_0$, stable under cobase change (whenever it exists).
\item[\hfill $C''_0 \ = $] set of arrows which are either in $C'_0$, or amalgamated sums of a \textit{nonempty} family of such arrows $A \xrightarrow{i_\alpha} B_\alpha$ with the same source $A$, $i_\alpha \in C'_0$; the arrow $A \longrightarrow B = \displaystyle\coprod_{\alpha}{}_{(A)} B_\alpha$ is in $C''_0$.
\end{itemize}

(Condition 1.16 implies that the elements of $C'_0$ are cosquarable, and this suffices, together with the condition of stability of $\cM$ under filtering $\varinjlim$, to ensure that the infinite amalgamated sums envisaged exist in $\cM$.)

\textbf{N.B.} $C''_0$ is \textit{equal} to the smallest set of arrows containing $C_0$ and stable under cobase change and under small amalgamated sums. \texttt{[If small sums exist in $\cM$, then $C''_0$ is also the set of arrows obtained by cobase change from sums of arrows of $C_0$.]}

\setcounter{footnote}{1}% the edition numbers these footnotes 2 and 3
\begin{itemize}[leftmargin=10em, labelwidth=9em, labelsep=1em, align=left]
\item[(1.19)\hfill $C'''_0 \ = $] smallest subset of $\Fl(\cM)$ containing $C''_0$, and moreover stable under strict ordinal direct limits, \textit{i.e.} for an ordinal direct system $(A_\alpha)_{\alpha<\alpha_0}$, such that for $\alpha$ a limit ordinal we have $A_\alpha \xleftarrow{\ \sim\ } \varinjlim\limits_{\alpha' \in I_\alpha} A_{\alpha'}$, if the transition morphisms $i_{\beta,\alpha} : A_\alpha \longrightarrow A_\beta$ are in $C'''_0$, so too the canonical morphisms
\[
A_\alpha \xrightarrow{\ i_\alpha\ } \varinjlim_{\beta \in I_{\alpha_0}} A_\beta \ . \quad \footnotemark
\]
\end{itemize}
\footnotetext{$C_0 \subset C'_0 \subset C''_0 \subset C'''_0$}% footnote 2: mark printed in the display, text printed at the foot of p. 8

\printedpage{8}
% head: A. Grothendieck : Dérivateurs, Ch. XV, Édition des Universités de Montpellier II et Paris VII
% foot: Version 5 novembre 2010
\msp{page 9}

Proposition 2, p. 4 of XIII ensures that we have
\[
C'''_0 \ \subset \ \widetilde{C_0} \ .
\]
This being so, here is the second essential property of the factorization $f = p_1(f)i_1(f)$ (the first being 1.17):

\begin{itemize}[leftmargin=5em, labelwidth=4em, labelsep=1em, align=left]
\item[(1.20)] \qquad\qquad For every arrow $f \in \Fl(\cM)$, we have $i_1(f) \in C''_0$.
\end{itemize}

It follows at once, by ordinal induction:

\begin{itemize}[leftmargin=5em, labelwidth=4em, labelsep=1em, align=left]
\item[(1.21)] \qquad For every arrow $f \in \Fl(\cM)$ and every ordinal $\alpha \in \fU$, we have $i_\alpha(f) \in C'''_0$, a fortiori $i_\alpha(f) \in \widetilde{C_0}$.
\end{itemize}

So if $\alpha$ is large relative to $\pi$ (\textit{cf.} cor. 2, page 7), the factorization $f = p_\alpha(f)i_\alpha(f)$ satisfies
\[
\left\{
\begin{array}{l}
f = p_\alpha(f)i_\alpha(f)\text{, with } i_\alpha(f) \in C'''_0 \subset \widetilde{C_0} = \Phi \text{ \texttt{[not to be}}\\
\text{\texttt{confused with }}\Phi_\alpha\text{\texttt{]}, } p_\alpha(f) \in C_{0*} = \Psi \text{, when } \alpha \text{ is}\\
\text{large relative to } \pi\text{, the sources of the } i \in C_0 \ \pi\text{-accessible.}
\end{array}
\right.
\tag*{(1.22)}
\]

It remains to construct the functor $f \longmapsto X_1(f)$, with $i_1(f)$, $p_1(f)$, satisfying conditions (1.17) and (1.20), under the hypothesis (1.16) that the $i \in C_0$ are cosquarable. We then consider, for $f \in \Fl(\cM)$ fixed, the set $\Delta(f)$ (or $\Delta(C_0, f)$) of \textit{all} the diagrams $D$ of the form (1.15), with $i \in C_0$, and we put
\[
X_1(f) \ = \ \coprod_{D \in \Delta(f)}{}_{(X)} X_D
\tag*{(1.23)}
\]
(amalgamated sum under $X$, \textit{i.e.} sum indexed
\msp{page 10}
by $\Delta(f)$ in $X\backslash\cM$), the morphism $X \xrightarrow{i_1(f)} X_1(f)$ being the costructural morphism, and the morphism $X_1(f) \xrightarrow{p_1(f)} Y$ being the one with the $p_D$ as components. In short, we proceed ``as economically as possible'' (and as naturally as possible) to ensure the validity of (1.17), in the equivalent form (1.15). On the other hand, by definition of $C'_0$, $C''_0$, the $i_D$ are in $C'_0$, and $i_1(f)$ is in $C''_0$, which is (1.20)\footnote{One should all the same make explicit the functorial dependence of $X_1(f)$ on $f$\,!!!}.

To sum up:

\textbf{Theorem 1.} \textit{Let $\cM$ be a $\fU$-category stable under filtering $\fU$-direct limits, $C_0 \subset \Fl(\cM)$ a small set of arrows of $\cM$. We assume the arrows $i \in C_0$ to be cosquarable. This ensures the definition of an endofunctor $f \longmapsto p_1(f)$ of $\Flu(\cM)$, and of a functorial homomorphism $i_1 : \mathrm{id}_{\Flu(\cM)} \longrightarrow p_1$,} i.e. \textit{$i_1(f) : f \longrightarrow p_1(f)$, inducing the identity on the target objects of $f$ and of $p_1(f)$,}
\[
\begin{tikzcd}[column sep=small]
X \arrow[rr, "i_1(f)"] \arrow[dr, "f"'] & & X_1(f) \arrow[dl, "p_1(f)"] \\
 & Y \ , &
\end{tikzcd}
\]